% Generated draft: requires mathematical and source review.
\documentclass[11pt]{article}
\usepackage[margin=25mm]{geometry}
\usepackage{fontspec}
\setmainfont{Latin Modern Roman}
\usepackage{amsmath,amssymb,mathtools,unicode-math}
\setmathfont{Latin Modern Math}
\usepackage{enumitem,needspace,xurl,hyperref}
\setlength{\parindent}{0pt}
\setlength{\parskip}{5pt}
\setlength{\emergencystretch}{3em}
\newcommand{\N}{\mathbb{N}}
\newcommand{\Z}{\mathbb{Z}}
\newcommand{\Q}{\mathbb{Q}}
\newcommand{\R}{\mathbb{R}}
\newcommand{\C}{\mathbb{C}}
\newcommand{\F}{\mathbb{F}}
\newcommand{\A}{\mathbb{A}}
\newcommand{\PP}{\mathbb{P}}
\newcommand{\E}{\mathbb{E}}
\newcommand{\eps}{\varepsilon}
\newcommand{\dd}{\,\mathrm{d}}
\begin{document}
\section*{UnsolvedMath ID: 6200077}
\textbf{Source problem code:} catalogue entry
\subsection{Definitions and mathematical statement}
Let $M$ be a closed hyperbolic $n$-manifold and $G=\pi_1(M)$. A group $\Gamma$ is residually finite if every $\gamma\ne1$ has nonidentity image under some homomorphism from $\Gamma$ to a finite group.

\textbf{Problem (Misha Kapovich).} For a short exact sequence
\[
 1\longrightarrow A\longrightarrow\Gamma\longrightarrow G\longrightarrow1,
 \qquad A\text{ printed as ``Zp''},
\]
is $\Gamma$ residually finite? The source then asks, as a reformulation, whether there is a finite-index subgroup $G'\leq G$ for which
\[
 \operatorname{res}:H^2(G,A)\longrightarrow H^2(G',A)
 \quad\text{is the zero map}. \tag{1}
\]
The kernel notation is damaged/ambiguous: $\Z/p\Z$, $\Z^p$, and $p$-adic integers are different possibilities, and no range for $p$ is supplied. Formula (1) must retain the intended coefficient meaning rather than silently choose one.

Under the plausible finite-cyclic interpretation $A=\Z/p\Z$, an untwisted $H^2(G,A)$ describes central extensions. The short exact sequence as printed does not explicitly assume centrality or specify the conjugation action on $A$; a nontrivial action instead uses the corresponding coefficient module. These are essential qualifications of the cohomological formulation.

For a fixed central extension with class $e\in H^2(G,A)$, restriction vanishing for that extension means $\operatorname{res}(e)=0$, i.e. splitting over $G'$. The source asks the stronger-looking uniform zero-map condition for all classes. In a finite-coefficient setting its relationship to the universal question over all extensions requires the relevant finiteness and passage-to-common-subgroup argument; it is not confused with a claim about one class without those conditions. The finite versus infinite kernel distinction also matters for any residual-finiteness/splitting equivalence.

The summary reports partial progress, attributes a ``$p=1$'' case to Agol, and mentions poorly specified residual-finiteness terminology without a theorem reference. That $p=1$ description does not resolve the ambiguous coefficient notation, and no precise dimensional or coefficient scope is authenticated. The general problem is reported open in the supplied evidence.
\subsection{Short English statement}
Are the stated extensions of closed hyperbolic-manifold groups residually finite, equivalently in the intended setting can finite-index restriction kill the relevant degree-two cohomology? The kernel notation Zp, centrality/module action, and single-class versus uniform restriction quantifiers need source confirmation.
\subsection{Sources}
\begin{itemize}
\item Local UnsolvedMath registry, permanent ID 6200077, source code catalogue entry, ``Boundaries of Groups and Kleinian Groups --- Problem 77''; proposer: Misha Kapovich, proposed year 2005. Status partially solved, classification PARTIAL-PROGRESS, updated 2026-08-06.
\item Kapovich, \emph{Problems on Boundaries of Groups and Kleinian Groups} (2005), Problem 77, PDF page 20, \url{https://www.math.ucdavis.edu/~kapovich/EPR/problems.pdf}. Reported PDF-text-normalized extraction and access: 2026-07-29. Current status NEEDS\_REVIEW; the source labels it a problem. Rights: public author-hosted PDF; statement-reuse terms NEEDS\_REVIEW.
\item The summary names Agol and a $p=1$ case but gives no primary reference or precise theorem; its additional ``jp'-RF''/linear-by-hyperbolic wording is not defined in the record.
\end{itemize}
\end{document}
